\documentclass{article}
\usepackage{amsmath, amssymb, amsthm}
\usepackage{hyperref}
\usepackage{geometry}
\geometry{margin=1in}

\title{Erd\H{o}s Problem \#850}
\date{Last updated: 2025-08-31}
\author{Source: erdosproblems.com}

\begin{document}
\maketitle

\noindent\textbf{Status:} OPEN \quad \textbf{No prize}

\noindent\textbf{Tags:} number theory, primes

\medskip
\noindent\textbf{Problem Statement:}

Can there exist two distinct integers $x$ and $y$ such that $x,y$ have the same prime factors, $x+1,y+1$ have the same prime factors, and $x+2,y+2$ also have the same prime factors?

\bigskip
\noindent\textbf{Remarks:}

This is sometimes known as the Erd\H{o}s-Woods conjecture.

For just $x,y$ and $x+1,y+1$ one can take\[x=2(2^r-1)\]and\[y = x(x+2).\]Erd\H{o}s also asked whether there are any other examples. Makowski \cite{Ma68} observed that $x=75$ and $y=1215$ is another example, since\[75 = 3\cdot 5^2 \textrm{ and }1215 = 3^5\cdot 5\]while\[76 = 2^2\cdot 19\textrm{ and }1216 = 2^6\cdot 19.\](This example was also found independently by Matthew Bolan, and by Dubickas, who posed it as part of the 2024 team selection test in Lithuania.) No other examples are known. This sequence is listed as A343101 at the OEIS.

Shorey and Tijdeman \cite{ShTi16} have shown that, assuming a strong form of the ABC conjecture due to Baker, then the answer to the original problem is no.

See also [677]. 

The case of $x,y$ and $x+1,y+1$ appeared as Problem 1 in the Third Benelux Mathematical Olympiad 2011. 

This problem is discussed in problem B19 of Guy's collection \cite{Gu04}.

\bigskip
\noindent\textbf{References:}

[Gu04] Guy, Richard K., Unsolved problems in number theory. (2004), xviii+437.
 
 [Ma68] Makowski, Andrzej, On a problem of {E}rd\H{o}s. Enseign. Math. (2) (1968), 193.
 
 [ShTi16] Shorey, Tarlok N. and Tijdeman, Rob, Arithmetic properties of blocks of consecutive integers. (2016), 455--471.

\bigskip
\noindent\small{Source: \url{https://www.erdosproblems.com/850}}
\end{document}
